\section{UK research stream A: Inheritance, inequality and an extension of the IFS work}
\label{sec:ifs}

\subsection{Starting point and contribution}
The most direct starting point is Bourquin, Joyce and Sturrock, \emph{Inheritances and Inequality within Generations} (IFS R173, 2020), with the follow-on \emph{Inheritances and Inequality over the Life Cycle} (IFS R188, 2021). The first combines information from the ONS Longitudinal Study, ELSA and WAS to project parental wealth, mortality and the size and age of future inheritances for people born in the 1960s--1980s. The second considers the implications for lifetime resources and the anticipation of future inheritances.

The central opportunity is to \textbf{move from projected inheritances to better measurement and validation against realised transfers}. This is a data contribution, independent of our separate theoretical work. IFS's 2020 construction makes three transparent but consequential simplifying assumptions: the first spouse leaves the estate to the surviving spouse; bequests pass to children when the second parent dies; and children divide the relevant estate equally. It also projects future estate values from observed older-household wealth patterns. These can be assessed rather than simply assumed.

A separate earlier IFS study, Crawford and Hood (2016), found that inheritance's measured effect on wealth inequality is sensitive to whether pension wealth is included. Any extension must distinguish marketable, housing, financial and pension wealth.

\subsection{Concrete research projects}
\begin{description}[style=nextline]
\item[A1. Update and validate inheritance-timing projections]
Reproduce published age-at-last-parent-death and projected age-at-inheritance distributions by birth cohort; compare with subsequent ELSA reported receipt years and newer survival histories. Report differences between death, probate and receipt rather than assuming they coincide. Parents still alive create right-censoring.
\item[A2. Test actual beneficiary allocation against equal division]
Use annotated modern wills to measure intended shares among children versus spouses, grandchildren, charities, trusts and non-family recipients. Where permitted, use \emph{completed estate accounts} to measure realised net shares, changes from the will and actual distribution dates. Report unequal allocation within families, including how it varies by estate value and number of potential heirs. Wills alone cannot establish the final allocation.
\item[A3. Update the projected distribution of inheritances]
Refresh parental wealth and asset composition with later WAS waves, survival assumptions with current ONS tables and the distribution of heirs using suitable census/family data. Produce cohort-specific projected amounts and explicit sensitivity to the actual timing and allocation patterns found in A1--A2. Compare projected and observed cohort aggregates where follow-up allows; do not claim individual prediction validation from distinct datasets.
\item[A4. Inheritance and wealth inequality at receipt and over time]
Using longitudinal WAS/ELSA observations, compare wealth positions before and after reported receipt, with sample weights and repeated wealth components. Report absolute transfer gaps, Gini/top shares and rank or percentile transitions; distinguish within-cohort inequality from between-family-background inequality. Estimate descriptive decomposition with and without pension wealth. A mechanical subtraction of inheritances from observed wealth is not a causal counterfactual.
\item[A5. Intergenerational persistence and inequality of opportunity]
Ask whether inherited transfers widen differences by parental wealth, education, housing tenure, region or family structure, including siblings and children with no inheritance. Separate the inequality \emph{of transfers} from the relationship between transfer receipt and pre-existing resources. Where parental wealth is not directly observed, use transparent proxies and sensitivity analysis.
\item[A6. Lifetime gifts plus bequests]
Add inter vivos gifts and loans to death-time inheritances to compare \emph{total lifetime transfers} with bequests alone. The timing of gifts may materially change lifetime opportunity even if the eventual inheritance is equal; avoid double counting gifts later reflected in a smaller estate.
\end{description}

\subsection{Required data, by priority}
\begin{description}[style=nextline]
\item[Feasible first: within licensed surveys]
ELSA respondent birth cohort, parental deaths, reported inheritance receipt year/value/source, wealth before and after; WAS transfer receipt, wealth components, household structure, parental-background measures if available, pensions and weights; ONS mortality tables. First perform exact variable and sample audits.
\item[High-value external validation: governed will/estate sample]
Deceased age and dates; legal family relationships; number of children; each named beneficiary and intended share; spouse and non-child bequests; estate value by asset type; final accounts identifying net distributions and payment dates. This can evaluate assumptions at the \emph{sample-distribution} level without linking named wills to anonymised survey participants.
\item[Ambitious administrative extension]
Custodian-approved family relationships, estate and recipient identities, transfer dates/values and wealth or tax outcomes, subject to confirmation that relevant data exist. Access to ONS LS and HMRC/ONS linkage routes is restricted and not guaranteed.
\end{description}

\subsection{Key empirical distinctions}
\begin{itemize}[nosep]
\item Absolute wealth gaps may rise even where relative inequality measures fall. Report both.
\item Transfers may reduce within-cohort relative inequality yet reinforce advantage between richer and poorer parental families.
\item Inheritances are selected events and expectations may change saving before receipt; pre/post differences alone do not identify a causal effect.
\item Probate grants, wills, taxable estates and surveys cover different populations; record selection and missingness.
\item Report non-recipients, partial recipients, first-parent inheritances, jointly held assets and gifts made before death.
\end{itemize}

\subsection{Proposed deliverables}
\begin{enumerate}[nosep]
\item A reproducible update of IFS-style cohort inheritance timing and amount projections, with a clear assumptions register.
\item A table comparing equal-split/second-parent assumptions against observed evidence from surveys and a governed will/estate sample.
\item A wealth-inequality decomposition using multiple wealth definitions, parental-background strata and longitudinal age profiles.
\item An explicit list of which results remain descriptive and which would require linked administrative data for causal or lifetime analysis.
\end{enumerate}

\subsection{Source studies}
\begin{itemize}[nosep]
\item Bourquin, Joyce and Sturrock (2020), IFS R173: \url{https://ifs.org.uk/publications/inheritances-and-inequality-within-generations}
\item Bourquin, Joyce and Sturrock (2021), IFS R188: \url{https://ifs.org.uk/publications/inheritances-and-inequality-over-life-cycle-what-will-they-mean-younger-generations}
\item Crawford and Hood (2016), inheritance and broad wealth measures: \url{https://ifs.org.uk/journals/lifetime-receipt-inheritances-and-distribution-wealth-england}
\end{itemize}
